\section{Why complete continuation matters}
\label{sec:limitations}

The source-cover criteria require bounds on complete raw recovery.
Neither disjoint exceptional spans nor small variation of raw reads
transported through the exceptional maps supplies such a bound by itself.
We give one native geometry that makes this distinction explicit.  It
has a quadratic fractional cover at one root and phase, even though all
exceptional-only transported profiles vanish on the nondistinguished
core families.  The example does not satisfy a fixed inverse-square
family-size envelope and is not a counterexample to an averaged cover
criterion.

We use the strictly clipped shadows of \cref{def:clipped-shadow}, with
deleted families acting as the identity. Their ordinary orbit includes
the seed and is not truncated at other paid vertices. It is therefore
distinct from the actual dependency cone after paid lower vertices have
been removed.

For a retained core $\mathcal K$, let $B_{\mathcal K}(u)$ be the raw
labels of its ordinary orbit from time $u$.  For a nonempty raw set $S$,
define its arbitrary-seed fractional core-cover number by
\begin{equation}
 \kappa_{\mathcal K}^{\rm all}(S)
 =\inf\left\{\sum_{u=0}^{T-1}w_u:
 w_u\ge0,\quad
 \sum_{u:\,a\in B_{\mathcal K}(u)}w_u\ge1
       \text{ for every }a\in S\right\}.
 \label{eq:arbitrary-seed-cover}
\end{equation}
Allowing every computation time as a seed can only make covering easier
than requiring the seeds to lie in the original joint orbit.

\begin{proposition}[Serial exceptional spans can require a quadratic cover]
\label{prop:serial-quadratic-cover}
\notformalized
For every integer $M\ge4$ and $b\ge1$, there is a prescribed native
geometry with
\[
 d=16M,\qquad L=d^2b,
\]
eight work heads, unit moves or pauses, at most three entries per work
cell, fewer than $13L$ time vertices, and $L$ distinct increasing output
times, with the following properties.  In particular the construction
allows the logarithmic entry width $b=b(d)$.

The complete actual return families consist of a core $\mathcal K_0$
whose full spans partition into two disjoint-span classes, and two
exceptional families $F,G$ with disjoint full spans.  There are no
residual families.  A period-$M$ deletion of complete core families,
together with payment of every selected outer lower endpoint, produces
a retained core $\mathcal K$, a paid set $\Gamma$, and an output root
$t$ such that, writing $S(t)$ for its complete ordinary raw orbit and
$Q=(M-1)^2$,
\begin{equation}
 \kappa_{\mathcal K}^{\rm all}(S(t))=Q,
 \qquad
 \gaparea(S(t))\le\frac1{2Q}+\frac1{2L}.
 \label{eq:serial-cover-conclusions}
\end{equation}
The same cover equality holds if covering seeds must lie in the joint
orbit of $t$.  The paid physical dependency cone at $t$ contains $Q$
raw coordinates with circular gaps at most $\lceil L/Q\rceil$.

There is one core family $H$ such that, on the entire full span of every
other original core family, $x\circ w$ is identically zero for every
finite word $w$ in the two exceptional shadows, including the empty
word.  Here $x$ is the raw input-head position.  Thus all these
exceptional-only variation tails vanish, despite the quadratic cover.
\end{proposition}

\begin{informalproof}
Put
\[
 m=M,\quad n=8m,\quad h=4nb,
 \quad L=nh=4n^2b,
 \quad A=4L,\quad E=11L,
 \quad T=12L+h+1.
\]
All heads begin at zero.  Four heads form two copies of the following
bank construction, beginning at times $A$ and $E$.  At time $D+2ih$,
$0\le i<n/2$, head $i\bmod2$ moves up $h$, down $h$, and up $h$ from
its current base, each leg at unit speed; it waits until its next bank.
Before and after these movements it remains parked.  The base increases
by $h$ after each bank on that head.  Direct inspection of the successive
visits gives exactly the families
\begin{equation}
 \begin{split}
 C_j^D:\quad &\text{span }(D+jh,D+(j+2)h],
       \quad\text{owner }D+(j+1)h,\\
 &\text{hops }(D+(j+1)h-u,D+(j+1)h+u),\quad1\le u\le h,
 \end{split}
 \label{eq:serial-bank-families}
\end{equation}
for $D\in\{A,E\}$ and $0\le j<n$.  The two families in each bank
come from its upper and lower turn respectively.  Its interior cells
are entered on its three legs; a later bank uses a fresh spatial
interval.  The common boundary cell also has at most three entries.
Parking adds only adjacent previous visits.  Thus
\cref{eq:serial-bank-families} is the complete nonadjacent family list
on those four heads.  Even and odd $j$ give two disjoint-span classes.

A fifth head supplies $F$: it waits through $7L$, rises to $L$ by
$8L$, holds through $11L$, and descends to zero by $12L$.  Its one
complete family has
\[
 \text{span }(7L,12L],\quad\text{owner }8L,
 \quad\text{hops }(7L+a,12L-a),\quad0\le a<L.
\]
A sixth head supplies the core family $H$: it rises to $L$ by time
$L$, holds through $2L$, and descends to zero by $3L$.  Its family is
\[
 \text{span }(0,3L],\quad\text{owner }L,
 \quad\text{hops }(a,3L-a),\quad0\le a<L.
\]
There is no later ascent on either head, hence no additional lower-turn
family.  Their plateaux contribute adjacent visits.

The eighth head will supply a core bridge $B$.  Set $i=4m-1$ and
\[
 r_k=8L-(i-k)h-h/2,\qquad0\le k<m.
\]
This head waits through $5L$, rises to $L$ by $6L$, descends at unit
speed to height $m$ by $7L-m$, and then waits.  It descends by one at
each time $r_0,\ldots,r_{m-1}$ and parks at zero.  These times increase,
and $r_0>7L-m$, so the schedule is legal.  Its complete family has
owner $6L$, span $(5L,r_{m-1}]$, and hops
\[
 (5L+a,7L-a)\quad(m\le a<L),\qquad
 (5L+m-1-k,r_k)\quad(0\le k<m).
\]
In particular $B$ has $L$ actual hops.  Among bank spans it overlaps
only $C_{n-1}^{A}$: the preceding span ends exactly at $5L$.  Assign
$B$ to the even class and $H$ to that class as well.  The entire core
still has two disjoint-span classes; the early overlap graph is the
path $C_0^A,\ldots,C_{n-1}^A,B$.

The input head follows
\begin{equation}
 x(t)=
 \begin{cases}
 t,&0\le t\le L-1,\\
 2L-2-t,&L-1\le t\le2L-2,\\
 0,&t\ge2L-2.
 \end{cases}
 \label{eq:serial-raw-walk}
\end{equation}
All raw coordinates have been visited before the prescribed output
times $E,E+1,\ldots,E+L-1$.

We next specify the seventh head and the paid phase.  On the early
overlap path give $C_j^A$ rank $\min\{\max\{j,1\},n-1\}$ and $B$
rank $n-1$.  On the late path give $C_j^E$ rank
$\min\{\max\{j,1\},n-2\}$; give $H$ rank one.  Delete rank residue
zero modulo $m$.  This deletes precisely
$C_m^D,C_{2m}^D,\ldots,C_{7m}^D$ in each bank component and retains
$H,B$.  These are also the path-spine ranks with the endpoint leaves
assigned their adjacent spine rank.

Choose
\[
 t=E+(4m-1)h+h/2.
\]
Its retained late-core orbit consists of
$v_k=E+(4m-1-k)h+h/2$, $0\le k\le m-2$.
Indeed each reflection reduces the window index by one, until the
next required family is $C_{3m}^E$, which was deleted.  The actual
$F$ hop from $v_k$ lands at $r_k$.  The actual $B$ hop then lands at
$5L+a_k$, where $a_k=m-1-k$.  The last early bank family sends this
point below the bridge's outer lower endpoint:
\[
 5L+a_k\longmapsto5L-a_k<5L.
\]
Keeping the complete older continuation gives the $Q=(m-1)^2$ times
\begin{equation}
 z_{k,\ell}=A+(n-\ell)h+
 \begin{cases}
 a_k,&\ell\text{ even},\\
 h-a_k,&\ell\text{ odd},
 \end{cases}
 \quad 0\le k\le m-2,\quad1\le\ell\le m-1.
 \label{eq:serial-older-grid}
\end{equation}
They are distinct: different $\ell$ give different windows, and within
a window the offsets $a_k$ are distinct.  They are all strictly
interior to their windows, their maximum is $5L-1$, and their lowest
window is $7m+1$.  The next older family $C_{7m}^A$ was deleted, so
these are the complete older bank continuations.

Order these times as $z_1<\cdots<z_Q$ and set
\[
 u_j=2L+\lfloor jL/Q\rfloor,\qquad0\le j<Q.
\]
This ordering defines the prescribed geometry; it is not an uncharged
sorting operation in a claimed algorithm.
The seventh head waits through $u_0=2L$.  It leaves cell $j$ at time
$u_j+1$, reaches cell $j+1$, and waits until its next specified move.
Its first peak is at $u_{Q-1}+1<3L$.  It holds at height $Q$ until
$z_1-1$, then descends by one at each $z_i$ and parks.  Since both
sequences are strictly increasing, all moves are unit moves or pauses.
Its only family $G$ has
\[
 \text{span }(2L,z_Q]=(2L,5L-1],\quad
 \text{owner }u_{Q-1}+1,\quad
 \text{hops }(u_{Q-i},z_i),\quad1\le i\le Q.
\]
The lower endpoints use the last outward visits, including all waits.
The full spans of $G$ and $F$ are disjoint and can share a third
disjoint-span class.  We have now listed all $2n+4$ complete families;
the new heads have no later ascent and no omitted opposite-side
family.  All work cells have at most three entries.

Pay every actual lower endpoint of the fourteen deleted bank families
and every outer lower endpoint in the full selected family set.  The
first set has $14h$ elements.  The second has $2n+4$ elements, of
which fourteen already belong to the first.  Therefore
\begin{equation}
 |\Gamma|=14h+2n-10.
 \label{eq:serial-paid-size}
\end{equation}
Every displayed path to a $z_i$ avoids these payments: all bank offsets
are interior, all its windows avoid the deleted lower windows, and
the $F$ and $B$ hops avoid their outer lowers.  For $1\le j<Q$ the
actual $G$ and $H$ hops then give
\[
 z_{Q-j}\longmapsto u_j
       \longmapsto L-\lfloor jL/Q\rfloor.
\]
These vertices also avoid $\Gamma$.  The last $G$ hop lands at the
paid outer lower $u_0=2L$ and is discarded in the physical cone.
Together with the root's direct read zero, the surviving paths prove
that the physical cone contains
\begin{equation}
 R=\{0\}\cup
   \{L-\lfloor jL/Q\rfloor:1\le j<Q\}.
 \label{eq:serial-physical-grid}
\end{equation}
All state predecessors, other work predecessors and direct reads are
still retained; they can only add sources to this witnessed subset.
Its $Q$ circular gaps are at most $\lceil L/Q\rceil$.  Inserting
points only decreases the sum of squared gaps, so
\[
 \gaparea(S(t))\le\gaparea(R)
       \le\frac{\lceil L/Q\rceil}{2L}
       \le\frac1{2Q}+\frac1{2L}.
\]

For the exact cover assertion, the ordinary joint time orbit consists
of six groups: the $m-1$ times $v_k$, their $m-1$ images $r_k$, the
$m-1$ images $5L+a_k$, the $Q$ times $z_i$, their $Q$ strict $G$
images, and the $Q$ subsequent $H$ images.  The calculations above
show that each is reached.  They also show that the union is closed
under all retained maps: within a bank the only change is the next
listed reflection; $B$ acts on the $r_k$; $F$ acts on the $v_k$;
$G$ acts on the $z_i$; and $H$ acts on the $G$ images.  Other maps
fix these points.  At the last $G$ return the ordinary clipping rule
replaces $2L$ by $2L+1$, and $H$ then sends it to $L-1$.  Thus
\begin{equation}
 S(t)=\{0,L-1\}\cup
      \{L-\lfloor jL/Q\rfloor:1\le j<Q\}.
 \label{eq:serial-complete-raw}
\end{equation}
The extra label $L-1$ here is an ordinary clipped source; its displayed
clipped path is not asserted to survive the physical payment.

Give each of the $Q$ nonzero labels in
\cref{eq:serial-complete-raw} weight one.  A complete retained core
orbit from any computation time contains at most one of them.  To see
this, after time $3L$ every core continuation stays after $3L$, where
$x=0$.  Before time $L$ the core orbit is a singleton.  Between $L$
and $2L$, the only possible extra core image is $H$'s held owner $L$,
whose raw value is $L-2$; this is not a target because $L/Q>256$.
Between $2L$ and $3L$, the root reads zero and the sole nonidentity
core map $H$ adds at most one raw value.  The endpoints obey the same
description.  Summing the fractional covering constraints for the
$Q$ target labels therefore gives
$\sum_u w_u\ge Q$.

Conversely, use the $Q$ seeds
\[
 2L+1,\quad u_1,\ldots,u_{Q-1}.
\]
Each lies in the joint orbit and has core raw set consisting of zero
and one different nonzero target in
\cref{eq:serial-complete-raw}.  Weight one on each seed covers $S(t)$.
This proves both asserted cover equalities.

Finally, $x$ is zero throughout $[2L,T-1]$.  Both exceptional shadows
preserve that interval: their nonidentity images stay strictly above
their outer lowers, which are $2L$ and $7L$, and elsewhere they are
the identity.  Every original core span except $H$ lies in that
interval.  Hence $x\circ w$ vanishes on each such full span for every
word $w$ of exceptional shadows.  This proves the profile assertion
and completes the construction.
\end{informalproof}

The obstruction in this example is the continuation after the bridge
exit.  Many distinct intermediate times remain distinct through the
older bank, and the final core return $H$ converts them into separated
raw coordinates.  Looking only at raw reads immediately after an
exceptional hop misses this last recovery step.  In particular,
\cref{prop:serial-quadratic-cover} rules out a universal pointwise
$O(M)$ fractional core-cover bound on the stated serial and zero-tail
premises.  It also rules out a uniform pointwise lower bound
$\gaparea(S(t))\ge c/(M+1)$ with $c>0$: along this family,
$(M+1)\gaparea(S(t))\to0$.

There are two essential restrictions on that conclusion.  First, the
example ties $d=16M$ to the growing phase parameter and supplies one
root and one phase.  It does not disprove a phase or column average
whose parameters are fixed in the order required by the uniform
source-cover theorem.  Second, selecting the four complete families
$\{H,B,F,G\}$ leaves maximum family size $h=L/n=o(L)$.
Thus \cref{cor:finite-skeleton} already excludes near-full correctness
on this sequence of geometries; it does not provide a new hard native
class.  Its $2n$ equal bank families also violate every common fixed
inverse-square envelope: the $n$th largest core family has at least
$h=L/n$ hops, so a bound $h_j\le CL/j^2$ would force $C\ge n$.
These are prescribed legal geometries and source sets, not a
transposition algorithm or a proof that every reachable source carries
necessary information.  An affordable cut, an averaged cover bound,
or a different complete recovery certificate may still handle them.
